\documentclass{article}
\usepackage{amsmath, amssymb, amsthm}
\usepackage{hyperref}
\usepackage{geometry}
\geometry{margin=1in}

\title{Erd\H{o}s Problem \#670}
\date{Last updated: 2026-04-16}
\author{Source: erdosproblems.com}

\begin{document}
\maketitle

\noindent\textbf{Status:} OPEN \quad \textbf{No prize}

\noindent\textbf{Tags:} geometry, distances

\medskip
\noindent\textbf{Problem Statement:}

Let $A\subseteq \mathbb{R}^d$ be a set of $n$ points such that all pairwise distances differ by at least $1$. Is the diameter of $A$ at least $(1+o(1))n^2$?

\bigskip
\noindent\textbf{Remarks:}

The lower bound of $\binom{n}{2}$ for the diameter is trivial. Erd\H{o}s \cite{Er97f} proved the claim when $d=1$. The quantifiers are slightly ambiguous, but likely Erd\H{o}s had in mind the regime where $d$ is fixed and the $o(1)$ term $\to 0$ as $n\to \infty$ (perhaps at a rate which depends on $d$).

This was disproved if $n$ grows with $d$ by Ho \cite{Ho26}: there exist infinitely many $n$ and, with $d=n^2-n$, a set $A\subset \mathbb{R}^d$ such that all pairwise distances differ by at least $1$, with diameter at most\[\left(1-\frac{1}{\pi^2}+o(1)\right)n^2\approx 0.898n^2.\]

\bigskip
\noindent\textbf{References:}

[Er97f] Erd\H{o}s, Paul, Some unsolved problems. Combinatorics, geometry and probability (Cambridge, 1993) (1997), 1-10.
 
 [Ho26] B. S. Ho, Erd\H{o}s's diameter conjecture for separated distances fails in high dimensions. arXiv:2604.15305 (2026).

\bigskip
\noindent\small{Source: \url{https://www.erdosproblems.com/670}}
\end{document}
